\section{Discussion}
\label{sec:conclusion}

We examine minority participation in online moral discussion through
reception, comment--reply exchanges, and reader responses. In
high-agreement AITA discussions, minority comments receive lower scores yet
attract more replies, and challenging replies frequently articulate moral
reasons absent from their parent comments. In the survey, readers in
\emph{Balanced} report greater consideration of alternative judgments and
lower mean judgment certainty than readers in \emph{Highly Scored}, while
observed majority-alignment proportions are similar. These findings
suggest that approval alone does not capture the role of dissent in
discussions of competing moral judgments.

Our contribution is to connect the social reception of dissent with the
reasoning articulated around it and the alternatives readers report
considering. Prior research links minority dissent to divergent thinking
and group cognitive complexity
\cite{nemeth1986differential,curseu2012socially}, and disagreement to the
ability to articulate opposing reasons \cite{price2002disagreement}.
We bring these concerns together through complementary analyses of public
exchanges and reader responses. Separating reply stance from added reasons
shows how rejection and reason-giving coexist within an exchange.
A minority judgment may contribute through responses that make competing
obligations and interpretations explicit, even when those responses
reject its conclusion. 
A broader range of articulated reasons, however, does not by itself
establish stronger arguments, respectful interaction, or better judgments.

Distinguishing blaming from defending dissent further brings the moral
direction of opposition into analyses of majority--minority relations.
Under high agreement, blaming dissent draws more extensive reply
engagement, while challenges to defending dissent have the highest
observed rate of additional reasons. Accounts of public blame as subject
to social demands for justification \cite{malle2022socially} offer a
possible interpretation of the reception asymmetry. Together, these
patterns show that the amount of interaction and the reasons added within
it capture different aspects of dissent. 

The observational and survey findings address complementary aspects of
engagement with dissent. The observational analysis identifies reasons
articulated in public exchanges, while the survey examines readers'
reported consideration under different comment displays. Together,
they motivate assessing consideration of alternatives alongside final
judgment alignment. They do not establish that the additional reasons
observed in replies explain the survey differences: the survey does not
isolate exposure to these reasons as a mechanism. Future studies could
trace which reasons readers encounter, recall, and use in explaining
their judgments.

The additional reasons articulated around dissent and readers'
reported consideration of alternatives motivate displays that make
both majority and minority judgments, and their supporting reasons,
accessible.
The findings motivate making both majority and minority judgments
accessible while also showing the moral reasons each invokes and those
developed in replies. Platforms could highlight these reasons to make
different grounds for interpreting the same situation explicit.
\citet{kriplean2012considerit} provide a design precedent by organizing
participation around shared pro and con reasons and reflection on tradeoffs.
Building on this approach, experiments could compare displays of opposing
threads alone with displays that also organize their reasons and
counterarguments. 
Drawing on work on argument repertoire
\citep{cappella2002argument,price2002disagreement}, evaluation could examine
the reasons readers articulate for their own and opposing judgments,
alongside judgment alignment and certainty. Delayed assessments could
examine which reasons readers retain, while follow-up discussions could
reveal how greater visibility of opposing views affects the tone of
interaction.

\runinheading{Limitations.}
We report several limitations of our study. Because blaming and defending
dissent occur in different posts, their observed differences may also
reflect the situations being discussed. 
Future work could examine whether these patterns persist across posts with comparable topics and conflict characteristics.
Although similar reception patterns appear in AITAH, both communities
center on judgments of interpersonal conflicts. Generalization to
communities with different topics, norms, and interaction structures
requires further validation.
The survey uses a limited set of scenarios and varies comment content, replies, and scores together.
Larger studies could cover more situations and vary content and social
cues separately. Open-ended and delayed responses could complement
readers' reported consideration of alternative reasons by showing which
reasons they articulate and retain.